\documentclass[10pt,a4paper]{amsart}
\usepackage[margin=19mm]{geometry}
\usepackage{amsmath,amssymb,mathtools}
\usepackage[scr=rsfs,cal=euler]{mathalfa}
\usepackage{xfrac}
\usepackage[unicode,hidelinks,bookmarks=false]{hyperref}
\newcommand{\ie}{i.e.\ }
\newcommand{\cf}{cf.\ }
\newcommand{\resp}{respectively\ }
\newcommand{\Subsection}{\subsection}
\providecommand{\nobreakdash}{\nobreak\hbox{-}}
\setlength{\emergencystretch}{2em}
\multlinegap0pt
\usepackage{amssymb}
\usepackage{mathtools}
\usepackage{algorithm}\usepackage{algpseudocode}
\usepackage[shortlabels]{enumitem}
\newtheorem{lemma}{Lemma}
\newtheorem{theorem}{Theorem}
\usepackage{cleveref}
\crefname{lemma}{Lemma}{Lemmas}
\crefname{theorem}{Theorem}{Theorems}
\newcommand{\N}{\mathbb{N}}
\title{On Effective Banach-Mazur Games and an application to the Poincar\'e Recurrence Theorem for Category}
\author{Prajval Koul, Satyadev Nandakumar}
\date{Source: arXiv:2506.11118v2 (CC BY 4.0)}
\begin{document}
\maketitle

\noindent\emph{Source and licence.} On Effective Banach-Mazur Games and an application to the Poincar\'e Recurrence Theorem for Category. Version arXiv:2506.11118v2 (CC BY 4.0), distributed by arXiv under the Creative Commons Attribution 4.0 International licence, http://creativecommons.org/licenses/by/4.0/, as recorded in the arXiv licence metadata for this exact version and shown on the abstract page https://arxiv.org/abs/2506.11118v2. Full licence notice: \texttt{licenses/arxiv-2506.11118-CC-BY-4.0.txt}.\par\smallskip

\noindent\textit{Editorial scope.} Counted unit: the article's theorem theorem (source0.tex lines 465--467). The article's complete original proof of it is delivered with it (source0.tex lines 468--510, one \texttt{proof} environment of 43 lines). No mathematics is written, repaired or re-derived: the statement and the proof are the article's own lines, in the article's own notation, and no retained line is substituted or re-flowed.  Retained uncounted as the source-local context the proof needs: lines 271--273; lines 327--329.  Nothing was omitted from any retained span, so the package carries no omission record.  No retained line contains a citation, so the wrapper carries no bibliography and no bibliography entry.  No retained line contains a figure environment, an image-inclusion command or drawing code, so no image is delivered and the package declares no asset dependency.  Editorial formatting only, in addition: an AMS \texttt{amsart} preamble that restates the article's own declarations and macros used by the retained lines, namely the environments \texttt{lemma}, \texttt{theorem} on separate counters, together with the standard packages those lines need.  The retained environments print in this document as Lemma 1, Theorem 1 in order of appearance, whereas the article prints the same items with the numbering of its own full document.  The complete native source of the article as distributed in the arXiv e-print for this exact version is bundled unchanged as \texttt{sources/179766/source0.tex}.
\begin{lemma}\label{lem:cedensecomp}
	The complement of a dense c.e. open set is effective nowhere dense. 
\end{lemma}

\begin{theorem}\label{thm:effbm1}
	In a strongly computable $T_0$ space $(X,\tau,\beta,\nu)$ with $M\coprod C=X$, the Banach-Mazur game $BM\langle M,C\rangle$ has an effective winning strategy for $P_2$ if, and only if, $M$ is an effective first category set in $X$. 
\end{theorem}

\begin{theorem}
	The set of non-Liouville numbers forms an effective first category set. 
\end{theorem}

\begin{proof}
The complement of the set of Liouville numbers can be written as
\begin{align*}
  E^c&=\mathbb{Q} \cup \bigcup_{n\in\N}
       \left(\bigcup_{\substack{p,q\in\mathbb{Z}\\q>1}}
  \left(\frac{p}{q}-\frac{1}{q^n},\frac{p}{q}+\frac{1}{q^n}\right)\right)^c.  
\end{align*} 
Observe that the set
$\bigcup_{\substack{p,q\in\mathbb{Z}\\q>1}}\left(\frac{p}{q}-\frac{1}{q^n},\frac{p}{q}+\frac{1}{q^n}\right)$ 
is dense in $\mathbb{R}$, since it contains all rationals. Moreover,
it is easily seen to be a c.e. open set. Hence, it is a dense c.e.
open set, and its complement is effective nowhere dense, by
\cref{lem:cedensecomp}.
%It should be noted that this set is meager. 

Now, consider the Banach-Mazur game $BM\langle E,E^c \rangle$.
Consider the following strategy for $P_2$. At stage $k$ of the game,
with $P_1$'s response being $G_{2k-1}\in\mathcal{G}$, $P_2$ plays the
set $G_{2k}\in\mathcal{G}$, where $G_{2k}$ is the first set in
$\mathcal{G}$ such that
\begin{align*}
  G_{2k}&\subseteq G_{2k-1}\mathbin{\big\backslash}\overline{\left(\bigcup_{\substack{p,q\in\mathbb{Z}\\q>1}}\left(\frac{p}{q}-\frac{1}{q^k},\frac{p}{q}+\frac{1}{q^k}\right)\right)^c}\\
	&=G_{2k-1}\cap\left(\overline{\left(\bigcup_{\substack{p,q\in\mathbb{Z}\\q>1}}\left(\frac{p}{q}-\frac{1}{q^k},\frac{p}{q}+\frac{1}{q^k}\right)\right)^c}\right)^c\\
	&=G_{2k-1}\cap\left(\bigcup_{\substack{p,q\in\mathbb{Z}\\q>1}}\left(\frac{p}{q}-\frac{1}{q^k},\frac{p}{q}+\frac{1}{q^k}\right)\right)^o\\
	&=G_{2k-1}\cap\bigcup_{\substack{p,q\in\mathbb{Z}\\q>1}}\left(\frac{p}{q}-\frac{1}{q^k},\frac{p}{q}+\frac{1}{q^k}\right)^o.
\end{align*}
Clearly
$\bigcup_{\substack{p,q\in\mathbb{Z}\\q>1}}\left(\frac{p}{q}-\frac{1}{q^n},\frac{p}{q}+\frac{1}{q^n}\right)^o$
is a non-empty c.e. open set in $\mathbb{R}$. Hence $G_{2k}$ is a
non-empty basic open set of the parent space. $P_2$ plays this set.

Now, observe that
\begin{align*}
  E^c\cap\bigcap G_k&=E^c\bigcap\left(\bigcap G_{2k-1}\cap G_{2k}\right)\\
                    &\subseteq E^c\bigcap\left(\bigcap_{k\in\N}G_{2k-1}\cap\left(G_{2k-1}\mathbin{\big\backslash}\overline{\left(\bigcup_{\substack{p,q\in\mathbb{Z}\\q>1}}\left(\frac{p}{q}-\frac{1}{q^k},\frac{p}{q}+\frac{1}{q^k}\right)\right)^c}\right)\right)\\
                    &=E^c\bigcap\left(\bigcap_{k\in\N}G_{2k-1}\mathbin{\big\backslash}\bigcup_{k\in\N}\overline{\left(\bigcup_{\substack{p,q\in\mathbb{Z}\\q>1}}\left(\frac{p}{q}-\frac{1}{q^k},\frac{p}{q}+\frac{1}{q^k}\right)\right)^c}\right)\\
                    &\subseteq E^c\bigcap\left(\bigcap_{k\in\N}G_{2k-1}\mathbin{\big\backslash}E^c\right)\\
                    &=\varnothing.
\end{align*}
Hence, the above strategy is an effective winning strategy for $P_2$
in the game $BM\langle E^c,\mathbb{R}\setminus E^c\rangle$. Therefore,
by \cref{thm:effbm1}, $E^c$ is an effective first category set.
\end{proof}

\end{document}
